% Chapter 5 - CrypTool Online Architecture
% Vorgezogen von ehemals Kapitel 6 (Rueckmeldung des Betreuers), damit alle
% Grundlagenkapitel (2-6) hintereinander stehen.
\section{CrypTool Online Architecture}
\label{chap:ctoArchitecture}

\subsection{Introduction to CrypTool}
\label{sec:introCryptool}
% TODO: siehe Strukturplan.tex, Abschnitt 5.1
% - CrypTool Projekt, CrypTool Online

\subsection{Architecture of CrypTool Online}
\label{sec:ctoArchitectureDetail}
% TODO: siehe Strukturplan.tex, Abschnitt 5.2
% - Workflow-Konzept, Plugins, Datenfluesse

\subsection{Existing Educational Components}
\label{sec:existingEducationalComponents}
% TODO: siehe Strukturplan.tex, Abschnitt 5.3
% - Analyse bestehender CTO-Plugins

\subsection{Related Work}
\label{sec:relatedWork}
% TODO: siehe Strukturplan.tex, Abschnitt 5.4
% - Vergleich mit bestehenden Lehrmitteln zu code-basierter/Post-Quanten-Kryptographie
% - Abgrenzung: was macht diese Arbeit anders/zusaetzlich?
